\documentclass[11pt]{article}

% ---------- Packages ----------
\usepackage[margin=1in]{geometry}
\usepackage{amsmath, amssymb, amsthm}
\usepackage{mathtools}
\usepackage{enumitem}
\usepackage{graphicx}
\usepackage{booktabs}
\usepackage{array}
\usepackage{hyperref}
\usepackage{xcolor}
\usepackage{fancyhdr}
\usepackage{lastpage}
\usepackage{listings}
\usepackage{courier}

% ---------- Hyperref ----------
\hypersetup{
    colorlinks=true,
    linkcolor=blue,
    urlcolor=blue
}

% ---------- Header/Footer ----------
\pagestyle{fancy}
\fancyhf{}

\lhead{MIT 6.1200J}
\chead{Problem Set 2}
\rhead{Nicolas Iskandar}

\cfoot{\thepage\ /\ \pageref{LastPage}}

% ---------- Custom Commands ----------
\newcommand{\R}{\mathbb{R}}
\newcommand{\N}{\mathbb{N}}
\newcommand{\Z}{\mathbb{Z}}

\newcommand{\answer}{
    \vspace{0.5em}
    \noindent\textbf{Answer.}
}

% ---------- Code Style ----------
\lstset{
    basicstyle=\ttfamily\small,
    breaklines=true,
    frame=single,
    columns=fullflexible,
    keepspaces=true
}

% ---------- Title ----------
\title{
    \vspace{-2em}
    MIT 6.1200J \\
    Problem Set 2
}
\author{Nicolas Iskandar}
\date{}

% ==================================================
\begin{document}

\maketitle

\section*{Problem 1}

\begin{enumerate}[label=\textbf{(\alph*)}]

\item \textbf{Theorem}: ``Every natural number is tall or wide, but not both.''

\answer

We must show that for every natural number $n$, $n$ is tall or wide, but not both.
Let $n$ be an arbitrary natural number. We proceed by cases.

\textbf{Case 1}: $n$ is tall.\\
We must show $n$ is not wide. [TODO: using the definition of tall, prove $n$ is not wide.]

\textbf{Case 2}: $n$ is not tall.\\
We must show $n$ is wide and not both. [TODO: using the definition of tall, prove $n$
is wide.] Since $n$ is wide and not tall, $n$ is not both tall and wide.

In both cases, $n$ is tall or wide, but not both. Since $n$ was arbitrary, every natural
number is tall or wide, but not both. $\blacksquare$

\item \textbf{Theorem}: ``There exists a happy natural number between 10 and 1000 that
has no happy proper divisors.''

\answer

We prove this existentially by choosing $n = 25$ as the witness.

We must verify three things:
\begin{enumerate}[label=(\roman*)]
    \item $25$ is between $10$ and $1000$: This is clear since $10 < 25 < 1000$. \checkmark
    \item $25$ is happy. [TODO: using the definition of happy, verify that $25$ is happy.]
    \item $25$ has no happy proper divisors. The proper divisors of $25$ are $1$ and $5$.
    [TODO: using the definition of happy, verify that neither $1$ nor $5$ is happy.]
\end{enumerate}

Therefore, $n = 25$ is a happy natural number between $10$ and $1000$ with no happy
proper divisors, which proves the theorem. $\blacksquare$

\item \textbf{Theorem}: ``There are no devious natural numbers.''

\answer

We prove the equivalent statement: for every natural number $n$, $n$ is not devious.

\textbf{Predicate}: Let $P(n)$ be the statement ``$n$ is not devious.''

We proceed by induction on $n$, with base cases $n = 0$ and $n = 1$.

\textbf{Base Case} ($n = 0$): We must show $P(0)$, i.e., that $0$ is not devious.
[TODO: using the definition of devious, prove $0$ is not devious.]

\textbf{Base Case} ($n = 1$): We must show $P(1)$, i.e., that $1$ is not devious.
[TODO: using the definition of devious, prove $1$ is not devious.]

\textbf{Inductive Step}: Let $n \geq 1$ be an arbitrary natural number. Assume $P(n)$
holds, i.e., that $n$ is not devious. We must show $P(n+1)$, i.e., that $n+1$ is not
devious.

\textit{Assumption}: $n$ is not devious.

\textit{Goal}: Show that $n + 1$ is not devious.

[TODO: using the definition of devious and the inductive hypothesis that $n$ is not
devious, prove $n+1$ is not devious.]

By induction, $P(n)$ holds for all $n \in \mathbb{N}$. Therefore, there are no devious
natural numbers. $\blacksquare$

\end{enumerate}

\section*{Problem 2}

\begin{enumerate}[label=\textbf{(\alph*)}]

\item Brynmor's mistake is in the inductive step. When substituting the inductive
hypothesis into $x_{i+1} = 6x_i - 8x_{i-1}$, he plugs in lower bounds for both $x_i$
and $x_{i-1}$. However, since $x_{i-1}$ appears with a \textbf{negative} coefficient
$(-8)$, replacing $x_{i-1}$ with its lower bound $2^{i-1}$ makes the expression
\textit{smaller}, not larger. That is, from $x_{i-1} \geq 2^{i-1}$ one can only
conclude $-8x_{i-1} \leq -8 \cdot 2^{i-1}$, so the inequality goes the wrong way:
\[
    x_{i+1} = 6x_i - 8x_{i-1} \geq 6 \cdot 2^i - 8x_{i-1}
\]
is valid, but substituting the lower bound of $x_{i-1}$ on the right gives an upper
bound, not a lower bound. Hence the step
\[
    x_{i+1} \geq 6 \cdot 2^i - 8 \cdot 2^{i-1}
\]
is \textbf{not justified}.

\item Erik's mistake is in the base case and the circular reasoning at the end. His
inductive hypothesis is $R(i) := ``x_i \geq 3x_{i-1}$,'' defined only for $i \geq 2$.
At the end, he concludes that since $x_1 > 0$ and $x_i \geq 3x_{i-1}$ for all
$i \geq 2$, we get $x_i > 0$ for all $i$. But this final step implicitly assumes
$x_{i-1} > 0$ in order to conclude $x_i \geq 3x_{i-1} > 0$, which is exactly what
he is trying to prove. That is, the conclusion $x_i > 0$ requires knowing
$x_{i-1} > 0$, making the argument \textbf{circular}. He never establishes by a
separate argument that positivity propagates; he simply assumes it.

\item Zach's mistake is in the inductive step when $i = 1$. He claims to prove
$S(i+1)$ using $S(i)$ and $S(i-1)$, but when $i = 1$, this requires $S(0)$, i.e.,
$x_0 = 4^0 = 1$. However, the sequence is only defined starting from $x_1$; $x_0$
does not exist. So the inductive step \textbf{fails for $i = 1$} because it invokes
an undefined term. Indeed, the claim $S(i) := ``x_i = 4^i$'' is false for $i = 2$
since $x_2 = 13 \neq 16 = 4^2$, confirming the proof is broken.

\item We prove that $x_i > 0$ for all $i \geq 1$.

Following Erik's approach, we use the stronger inductive hypothesis
\[
    P(i) := ``x_i \geq 3x_{i-1} \text{ and } x_{i-1} > 0\text{''}
\]
for $i \geq 2$, and handle $x_1 > 0$ separately. Equivalently, we prove by strong
induction that $x_i > 0$ for all $i \geq 1$ and $x_i \geq 3x_{i-1}$ for all
$i \geq 2$.

More precisely, let $P(i) := ``x_i \geq 3x_{i-1}$'' for $i \geq 2$, and we will
carry the positivity of all previous terms alongside.

\textbf{Claim}: For all $i \geq 1$, $x_i > 0$. Moreover, for all $i \geq 2$,
$x_i \geq 3x_{i-1}$.

\textbf{Proof} by strong induction.

\textbf{Base Cases}:
\begin{itemize}
    \item $x_1 = 4 > 0$. \checkmark
    \item $x_2 = 13 > 0$, and $x_2 = 13 \geq 12 = 3 \cdot 4 = 3x_1$. \checkmark
\end{itemize}

\textbf{Inductive Step}: Let $i \geq 2$. Assume that for all $1 \leq k \leq i$,
$x_k > 0$, and for all $2 \leq k \leq i$, $x_k \geq 3x_{k-1}$. We wish to show that
$x_{i+1} > 0$ and $x_{i+1} \geq 3x_i$.

Consider:
\[
    x_{i+1} = 6x_i - 8x_{i-1}.
\]
By the inductive hypothesis, $x_i \geq 3x_{i-1}$, so $x_{i-1} \leq \frac{x_i}{3}$,
and thus $8x_{i-1} \leq \frac{8}{3}x_i$. Therefore:
\[
    x_{i+1} = 6x_i - 8x_{i-1} \geq 6x_i - \frac{8}{3}x_i = \frac{10}{3}x_i.
\]
Since $x_i > 0$ by the inductive hypothesis, we have:
\[
    x_{i+1} \geq \frac{10}{3}x_i > 3x_i > 0.
\]
This establishes both $x_{i+1} \geq 3x_i$ and $x_{i+1} > 0$, completing the
inductive step.

By strong induction, $x_i > 0$ for all $i \geq 1$. $\blacksquare$

\end{enumerate}

\section*{Problem 3}

\begin{enumerate}[label=\textbf{(\alph*)}]

\item We compute the sequence of transitions from starting state $(5, 22, 0)$:

\[
\begin{array}{|c|c|c|c|}
\hline
r & s & a & \text{Rule applied} \\
\hline
5  & 22 & 0   & s \text{ even: } (2r,\, s/2,\, a) \\
10 & 11 & 0   & s \text{ odd: } (2r,\, (s-1)/2,\, a+r) \\
20 & 5  & 10  & s \text{ odd: } (2r,\, (s-1)/2,\, a+r) \\
40 & 2  & 30  & s \text{ even: } (2r,\, s/2,\, a) \\
80 & 1  & 30  & s \text{ odd: } (2r,\, (s-1)/2,\, a+r) \\
160& 0  & 110 & s = 0, \text{ no transition (final state)} \\
\hline
\end{array}
\]

The final state is $(160, 0, 110)$, and indeed $a_f = 110 = 5 \cdot 22$. \checkmark

\item We prove that $P(r, s, a) := ``rs + a = xy$'' is an invariant of the state machine.

\textbf{Base case}: The start state is $(x, y, 0)$. We verify:
\[
    r \cdot s + a = x \cdot y + 0 = xy. \checkmark
\]

\textbf{Inductive step}: Suppose $(r, s, a)$ is a reachable state satisfying $P$,
i.e., $rs + a = xy$, and $s > 0$ so a transition is possible. We consider two cases.

\textbf{Case 1}: $s$ is even. The next state is $(2r,\, s/2,\, a)$. We verify $P$:
\[
    (2r)\cdot\frac{s}{2} + a = rs + a = xy. \checkmark
\]

\textbf{Case 2}: $s$ is odd. The next state is $(2r,\, (s-1)/2,\, a+r)$.
We verify $P$:
\[
    (2r)\cdot\frac{s-1}{2} + (a + r) = r(s-1) + a + r = rs - r + a + r = rs + a = xy.
    \checkmark
\]

In both cases, $P$ holds for the next state. Therefore $P$ is an invariant. $\blacksquare$

\item A final state is one from which no transition is defined. By the problem
statement, transitions are defined precisely when $s > 0$. Therefore, the final states
are all triples of the form
\[
    (r, 0, a) \quad \text{for any } r, a \in \mathbb{N}.
\]

\item Let $(r_f, s_f, a_f)$ be any reachable final state. By part (c), $s_f = 0$.
By part (b), since $P$ is an invariant and the start state satisfies $P$, every
reachable state satisfies $P$. In particular:
\[
    r_f \cdot s_f + a_f = xy.
\]
Substituting $s_f = 0$:
\[
    r_f \cdot 0 + a_f = xy \implies a_f = xy.
\]
Therefore, in any reachable final state, $a_f = xy$. $\blacksquare$

\item We prove that $\text{bitcount}(s)$ is a strictly decreasing derived variable,
i.e., that every transition strictly decreases $\text{bitcount}(s)$.

Consider any transition from state $(r, s, a)$ with $s > 0$.

\textbf{Case 1}: $s$ is even. The new value of $s$ is $s/2$. Since $s > 0$, we have
$s \geq 2$, so $s/2 \geq 1$. Now:
\[
    \text{bitcount}(s/2) = \left\lceil \log_2\!\left(\frac{s}{2} + 1\right) \right\rceil
    < \left\lceil \log_2(s + 1) \right\rceil = \text{bitcount}(s),
\]
where the inequality holds because $\frac{s}{2} + 1 \leq s + 1$ for $s \geq 2$, so
$\log_2(\frac{s}{2}+1) < \log_2(s+1)$. More directly: dividing $s$ by $2$ removes
exactly the least significant bit of its binary representation, strictly reducing the
number of bits by at least 1.

\textbf{Case 2}: $s$ is odd. The new value of $s$ is $(s-1)/2$. Since $s \geq 1$
is odd, we have $s - 1 \geq 0$ is even, so $(s-1)/2 \in \mathbb{N}$. We verify that
the bit count strictly decreases:
\[
\begin{aligned}
\text{bitcount}\!\left(\frac{s-1}{2}\right)
&= \left\lceil \log_2\!\left(\frac{s-1}{2}+1\right) \right\rceil \\
&= \left\lceil \log_2\!\left(\frac{s+1}{2}\right) \right\rceil \\
&= \left\lceil \log_2(s+1)-1 \right\rceil \\
&= \left\lceil \log_2(s+1) \right\rceil - 1 \\
&= \text{bitcount}(s)-1.
\end{aligned}
\]
So $\text{bitcount}(s)$ decreases by exactly 1. \checkmark \\
Again, subtracting 1 from an odd number clears the last bit and then dividing by 2
shifts right, so the bit count strictly decreases.

In both cases, $\text{bitcount}(s)$ strictly decreases with every transition.
Since $\text{bitcount}(s) \in \mathbb{N}$ (it is a nonneg integer), a strictly
decreasing sequence of nonneg integers cannot continue forever. Therefore the algorithm
always terminates. $\blacksquare$

\item From part (e), $\text{bitcount}(s)$ decreases by exactly 1 at every step.
Starting from $(x, y, 0)$, the initial value of $s$ is $y$, so the derived variable
starts at $\text{bitcount}(y)$ and decreases by 1 at each step until it reaches $0$,
at which point $s = 0$ and the algorithm halts. Therefore the algorithm terminates
after exactly
\[
    \text{bitcount}(y) = \lceil \log_2(y+1) \rceil
\]
steps. $\blacksquare$

\end{enumerate}

\section*{Problem 4}

\textbf{Theorem.} \textit{If fewer than $n$ students among those in an $n \times n$
arrangement are initially enlightened, and additional students reach enlightenment only
by being taught by two or more neighbors, then there will be at least one student who
never reaches enlightenment.}

\textbf{Proof.}

We use a monotonically non-decreasing derived variable argument. Consider the
\textbf{perimeter} of the enlightened region, defined as the total number of edges
shared between an enlightened cell and an unenlightened cell (or the boundary of the
grid). More precisely, let
\[
    L(t) := \text{the number of edges on the boundary of the set of enlightened
    students on day } t.
\]

\textbf{Key Claim}: $L(t)$ is non-increasing over time, i.e., $L(t+1) \leq L(t)$
for every day $t$.

\textbf{Proof of Key Claim}: Consider a single student who becomes newly enlightened
on day $t+1$. This student must have had at least 2 enlightened neighbors on day $t$.
We analyze how the perimeter changes when this student becomes enlightened. Suppose
the newly enlightened student has $k \geq 2$ enlightened neighbors. Then:
\begin{itemize}
    \item The $k$ edges shared with enlightened neighbors are \textit{removed} from
    the perimeter (they no longer separate an enlightened cell from an unenlightened
    one).
    \item The $4 - k$ edges shared with unenlightened neighbors (or the boundary) are
    \textit{added} to the perimeter.
\end{itemize}
So the net change in perimeter from this one student is $(4 - k) - k = 4 - 2k$.
Since $k \geq 2$, we have $4 - 2k \leq 0$. Therefore each newly enlightened student
does not increase the perimeter, and so $L(t+1) \leq L(t)$. $\square$

\textbf{Bounding the initial perimeter}: Each enlightened student is a unit square,
contributing at most 4 edges to the perimeter. Therefore, if there are fewer than $n$
initially enlightened students:
\[
    L(0) \leq 4 \cdot (\text{number of initially enlightened students}) < 4n.
\]

\textbf{Bounding the final perimeter}: If \textit{all} $n^2$ students were
enlightened, the enlightened region would be the entire $n \times n$ grid, whose
perimeter is exactly $4n$:
\[
    L(\text{all enlightened}) = 4n.
\]

\textbf{Conclusion}: Since $L(t)$ is non-increasing and $L(0) < 4n$, we have
$L(t) < 4n$ for all $t$. But if every student were enlightened, the perimeter would
equal $4n$. Since the perimeter can never reach $4n$, the set of enlightened students
can never be the entire $n \times n$ grid. Therefore, at least one student never
reaches enlightenment. $\blacksquare$

\end{document}